\documentclass[a4paper,12pt]{article}
\usepackage{color}
\usepackage{graphicx, gensymb}
\usepackage{amsfonts, cancel, amssymb}
\usepackage{amsmath, array, esvect, esint}
\usepackage{typearea}
\usepackage[a4paper,left=2cm,right=2cm,top=2cm,bottom=3cm,bindingoffset=5mm]{geometry}
\newcommand\numberthis{\addtocounter{equation}{1}\tag{\theequation}}
\newcommand{\kla}{}

\begin{document}

\title{01 Atomradien}
\author{Joachim Schwardt}
\date{\today}
\maketitle

\section*{Aufgabe 1: Atomradius}
Abschätzung des Radius eines Heliumatoms über Atommasse $m=4,003\, u$, Dichte $\varrho = 0,13\, g/cm^3$ und räumlichen Füllfaktor $\alpha = 74\%$:
$$r=\sqrt[3]{\frac{3(V/\alpha)}{4\pi}}=\sqrt[3]{\frac{3m}{4\pi\varrho\alpha}}= 2.545\cdot 10^{-10}\, m\equiv 2,545\, \mathring{A}$$
Bestimmung des Bohr-Radius über $p\approx \hbar/r$:
\begin{align*}
E(r)&=T+V=\frac{1}{2m_e}\frac{\hbar^2}{r^2}-\frac{e^2}{4\pi\epsilon_0r}\\
0&\overset {!}{=}E'(a_0)=-\frac{\hbar^2}{m_ea_0^3}+\frac{e^2}{4\pi\epsilon_0a_0^2}\ \ \Rightarrow\ \ a_0 = \frac{4\pi\epsilon_0\hbar^2}{m_ee^2}\\
&\Rightarrow\ \ E(a_0)=\frac{m_ee^2}{4\pi\epsilon_0\hbar^2}\kla\left( {\frac{m_ee^2}{2m_e\cdot 4\pi\epsilon_0}-\frac{e^2}{4\pi\epsilon_0}} \right) =-\frac{m_ee^4}{32\pi^2\epsilon_0^2\hbar^2}
\end{align*}

\section*{Aufgabe 2: Kernradius}
Bestimmung des Abstandes $r_0$ der größten Annäherung für ein Proton mit einer kinetischen Energie von $T_0= 1\, MeV$, das frontal auf einen Goldkern trifft. Bestimmen Sie die kinetische Energie, wenn das Proton den Kernradius $R = 1.3A^{\frac{1}{3}}\cdot 10^{-15}m$ unterschreitet ($A\ -$ Massenzahl):
\begin{align*}
T(r)&=T_0+V=T_0-\frac{Ze^2}{4\pi\epsilon_0 r}\overset{!}{=}0\ \ \Rightarrow\ \ r_0 = \frac{Ze^2}{4\pi\epsilon_0 T_0}=1,14\cdot 10^{-13}\, m\equiv 15\, R\\
T_0^{(R)}&\overset{!}{=}\frac{Ze^2}{4\pi\epsilon_0 R}\equiv 15\, MeV
\end{align*}
\section*{Aufgabe 3: Grundlagen des BOHRschen Atommodells}
Der Virialsatz liefert:
\begin{align*}
T&=-\frac{1}{2}V\ \ \ \ \mathrm{und }\ \ \ \ E=\frac{1}{2}V=-E_{ion}\\
&\Rightarrow\ \ T=E_{ion}\\
L&=|\vv{L}|=|\vv{r}\times\vv{p}|\overset{\substack{\vv{r}\bot \vv{v}\\|}}{=}mrv=m\frac{-e^2}{4\pi\epsilon_0 V}\sqrt{\frac{2T}{m}}=\sqrt{2mE_{ion}}\frac{e^2}{4\pi\epsilon_0 \cdot 2E_{ion}}= \\
&=\frac{e^2}{4\pi\epsilon_0}\sqrt{\frac{m}{2E_{ion}}}=1,055\cdot 10^{-34}\, Js\equiv\hbar
\end{align*}
Mit der Quantisierungsbedingung von BOHR ($L_n=n\hbar$) folgt mit der Kernladungszahl $Z$ für die Energieniveaus $E_n$ von Wasserstoff bzw. Helium, sowie der reduzierten Masse $\mu = \frac{m_em_K}{m_e+m_K}$:
\begin{align*}
L&=\frac{Ze^2}{4\pi\epsilon_0}\sqrt{\frac{\mu}{-2E_n}}=n\hbar\\ &\Rightarrow\ \ E_n=-\frac{Z^2e^4}{32\pi^2\epsilon_0^2\hbar^2}\frac{\mu}{n^2}=-Z^2\underbrace{\left(\frac{e^2}{4\pi\epsilon_0}\right)^2\frac{m_e}{2\hbar^2}}_{=E_R}\frac{\mu}{m_e}\frac{1}{n^2}=\frac{-E_RZ^2\mu}{m_en^2}=-\frac{E_RZ^2}{n^2}\frac{1}{1+\frac{m_e}{m_K}}
\end{align*}
\section*{Z1: Übergänge im Deuteron und Tritium}
Mit der Formel für die Energieniveaus $E_n$ aus (3) erhalten wir für die $\beta :=\frac{m_K}{m_e}$ mit $Z=1$:
\begin{align*}
hf&=\frac{hc}{\lambda}=E_2-E_1=\frac{E_R}{1+\frac{1}{\beta}}\left(\frac{1}{1^2}-\frac{1}{2^2}\right)\ \ \Rightarrow\ \ 1+\frac{1}{\beta} = \frac{3E_R\lambda}{4hc}\\
&\Rightarrow\ \ \beta = \frac{4hc}{3E_R\lambda -4hc}=\frac{1}{\frac{3E_R\lambda}{4hc}-1}
\end{align*}
Relativverfahren mit $m_K=\frac{m_e}{\frac{m_e}{\mu}-1}$:
\begin{align*}
\frac{hc}{\lambda}&\propto\mu \ \ \Rightarrow\ \ \frac{\lambda_H}{\lambda_D}=\frac{\mu_D}{\mu_H}\ \ \Rightarrow\ \ m_D=\frac{m_e}{\frac{\lambda_D m_e}{\lambda_H\mu_H}-1}=\frac{m_e}{\frac{\lambda_D}{\lambda_H}(1+\frac{m_e}{m_H})-1}
\end{align*}



\end{document}